\subsection{Modified Wilson Mirror}

\paragraph{Overview:}
Constructed from transistors and resistors, a
Modified Wilson Mirror is an enhanced version of
the Wilson Current Mirror. It offers improved
performance in terms of output resistance and
accuracy of current replication.

\paragraph{Schematic:}
\begin{center}
  \begin{tikzpicture}
  	% Paths, nodes and wires:
  	\node[shape=rectangle, draw={rgb,255:red,255;green,0;blue,0}, line width=1pt, dash pattern={on 4pt off 2pt on 1pt off 2pt on 1pt off 2pt}, minimum width=1.965cm, minimum height=2.215cm] at (15.02, 2.688){};
  	\node[shape=rectangle, draw={rgb,255:red,255;green,0;blue,0}, line width=1pt, dash pattern={on 4pt off 2pt on 1pt off 2pt on 1pt off 2pt}, minimum width=1.965cm, minimum height=2.215cm] at (6, 7.145){};
  	\draw (3.23, 8.04) to[american voltage source, l_={$V_{cc}$}] (1.98, 8.04);
  	\node[nmos, arrowmos, xscale=-1](N1) at (3.5, 2.75){} node[anchor=east] at (N1.text){$M_0$};
  	\node[nmos, arrowmos](N2) at (6, 2.75){} node[anchor=west] at (N2.text){$M_1$};
  	\draw (3.5, 7.54) to[european resistor, l={$G_4$}] (3.5, 6.04);
  	\draw (6, 5.79) |- (6, 5.25);
  	\draw (4.75, 2.75) |- (6, 3.71);
  	\draw (5.02, 2.75) -- (4.48, 2.75);
  	\draw (3.5, 1.98) -| (3.5, 1.75) -| (6, 1.98);
  	\draw (4, 1) to[american voltage source, l={$V_{ee}$}] (2.75, 1);
  	\draw (4, 1) -| (4.75, 1.75);
  	\node[ground, rotate=-90] at (1.98, 8.04){};
  	\node[ground, rotate=-90] at (2.75, 1){};
  	\draw (6, 5.79) -- (6.5, 5.79);
  	\node[shape=rectangle, minimum width=0.965cm, minimum height=0.465cm](N3) at (7, 5.75){} node[anchor=center] at (N3.text){$V_{out}$};
  	\draw (6, 7.5) to[european resistor, l={\textcolor{rgb,255:red,255;green,0;blue,0}{$G_5$}}] (6, 6);
  	\draw (6, 5.79) -- (6, 6.02);
  	\draw (3.5, 7.54) |- (3.23, 8.04);
  	\path[draw={rgb,255:red,255;green,0;blue,0}] (3.5, 8.04) -| (6, 7.5);
  	\node[pmos, arrowmos](N4) at (15, 7.168){} node[anchor=west] at (N4.text){$M_1$};
  	\node[pmos, arrowmos, xscale=-1](N5) at (12.5, 7.168){} node[anchor=east] at (N5.text){$M_0$};
  	\draw (13.98, 7.168) -- (13.48, 7.168);
  	\draw (15, 4.71) -| (15, 4.27);
  	\draw (13.71, 7.168) |- (14.77, 6.25);
  	\draw (12.5, 3.75) to[european resistor, l={$G_4$}] (12.5, 2.25);
  	\draw (15, 3.75) to[european resistor, l={\textcolor{rgb,255:red,255;green,0;blue,0}{$G_5$}}] (15, 2.25);
  	\draw (15, 4.27) |- (15, 3.75);
  	\draw (12.5, 7.938) |- (15, 8.188) -| (15, 7.938);
  	\draw (15, 4.02) -- (15.5, 4.02);
  	\node[shape=rectangle, minimum width=0.965cm, minimum height=0.465cm](N6) at (16, 4.02){} node[anchor=center] at (N6.text){$V_{out}$};
  	\draw (12.5, 2.25) -- (12.5, 2) -| (13.75, 2);
  	\draw (12.5, 8.188) to[american voltage source, l_={$V_{cc}$}] (11.25, 8.188);
  	\node[ground, rotate=-90] at (11.25, 8.188){};
  	\draw (13, 1.25) to[american voltage source, l={$V_{ee}$}] (11.75, 1.25);
  	\draw (13, 1.25) -| (13.75, 2);
  	\node[ground, rotate=-90] at (11.75, 1.25){};
  	\draw (10.25, 10.75) -- (10.25, -0);
  	\draw (1, 9.25) -- (19, 9.25);
  	\node[shape=rectangle, minimum width=3.84cm, minimum height=1.09cm] at (4.813, 10.063){} node[anchor=north, align=center, text width=3.452cm, inner sep=6pt] at (4.813, 10.625){Full Schematic\\(Sink)};
  	\node[shape=rectangle, minimum width=3.84cm, minimum height=1.09cm] at (13.562, 10.063){} node[anchor=north, align=center, text width=3.452cm, inner sep=6pt] at (13.562, 10.625){Full Schematic\\(Source)};
  	\node[shape=rectangle, minimum width=3.09cm, minimum height=2.465cm] at (8.687, 7){} node[anchor=north west, align=left, text width=2.702cm, inner sep=6pt] at (7.125, 8.25){\textcolor{rgb,255:red,255;green,0;blue,0}{\footnotesize Exists Only For Analysis. In  Implementation take output directly from V\_\{out\}}};
  	\path[draw={rgb,255:red,255;green,0;blue,0}] (13.75, 2) -| (15, 2.25);
  	\node[shape=rectangle, minimum width=3.09cm, minimum height=2.465cm] at (17.708, 2.562){} node[anchor=north west, align=left, text width=2.702cm, inner sep=6pt] at (16.145, 3.812){\textcolor{rgb,255:red,255;green,0;blue,0}{\footnotesize Exists Only For Analysis. In  Implementation take output directly from V\_\{out\}}};
  	\node[shape=rectangle, minimum width=0.965cm, minimum height=0.465cm](N7) at (7, 3.71){} node[anchor=center] at (N7.text){$V_{x}$};
  	\draw (6, 3.71) |- (6.5, 3.71);
  	\node[shape=rectangle, minimum width=0.965cm, minimum height=0.465cm](N8) at (16, 6.25){} node[anchor=center] at (N8.text){$V_{x}$};
  	\draw (14.98, 6.25) |- (15.48, 6.25);
  	\node[shape=rectangle, minimum width=0.965cm, minimum height=0.465cm](N9) at (4.75, 3.96){} node[anchor=center] at (N9.text){$I_{G_1}$};
  	\node[shape=rectangle, minimum width=0.965cm, minimum height=0.465cm](N10) at (13.71, 6){} node[anchor=center] at (N10.text){$I_{G_1}$};
  	\draw (12.5, 4.562) -- (12.5, 3.75);
  	\draw (1, -0) -- (19, -0);
  	\draw (12.5, -7.063) to[american current source, l_={$g_{m_1}v_{gs}$}] (12.5, -8.563);
  	\draw (14, -7.063) to[european resistor, l_={$g_{o_1}$}] (14, -8.563);
  	\draw (12.5, -8.563) -- (14, -8.563);
  	\draw (14, -7.063) -- (12.5, -7.063);
  	\draw (12.5, -8.563) -- (11.25, -8.563);
  	\draw (11.25, -7.063) -- (10.5, -7.063);
  	\node[shape=rectangle, minimum width=0.465cm, minimum height=0.465cm] at (11.25, -7.313){} node[anchor=center, align=center, text width=0.077cm, inner sep=6pt] at (11.25, -7.313){\footnotesize G};
  	\node[shape=rectangle, minimum width=0.465cm, minimum height=0.465cm] at (11.25, -8.313){} node[anchor=center, align=center, text width=0.077cm, inner sep=6pt] at (11.25, -8.313){\footnotesize S};
  	\node[shape=rectangle, minimum width=0.465cm, minimum height=0.465cm] at (14.5, -7.313){} node[anchor=center, align=center, text width=0.077cm, inner sep=6pt] at (14.5, -7.313){\footnotesize D};
  	\draw (10.5, -7.063) -- (9.75, -7.063);
  	\draw (8.25, -7.063) to[american current source, l={$g_{m_0}v_{gs}$}] (8.25, -8.563);
  	\draw (6.75, -7.063) to[european resistor, l={$g_{o_0}$}] (6.75, -8.563);
  	\draw (8.25, -7.063) -- (6.75, -7.063);
  	\draw (8.25, -8.563) -- (6.75, -8.563);
  	\draw (6.75, -7.063) -| (6.75, -5.563);
  	\draw (9.75, -8.563) -- (8.25, -8.563);
  	\node[shape=rectangle, minimum width=0.465cm, minimum height=0.465cm] at (9.75, -7.313){} node[anchor=center, align=center, text width=0.077cm, inner sep=6pt] at (9.75, -7.313){\footnotesize G};
  	\node[shape=rectangle, minimum width=0.465cm, minimum height=0.465cm] at (9.75, -8.313){} node[anchor=center, align=center, text width=0.077cm, inner sep=6pt] at (9.75, -8.313){\footnotesize S};
  	\node[shape=rectangle, minimum width=0.465cm, minimum height=0.465cm] at (6.25, -7.313){} node[anchor=center, align=center, text width=0.077cm, inner sep=6pt] at (6.25, -7.313){\footnotesize D};
  	\draw (6.75, -1.25) to[european resistor, l={$G_4$}] (6.75, -2.75);
  	\draw (6.75, -2.75) -| (6.75, -3.25);
  	\draw (6.75, -1.25) -| (6.75, -0.75) -- (7.5, -0.75) -| (7.5, -1);
  	\draw (14, -2.063) to[european resistor, l={$G_5$}] (14, -3.563);
  	\draw (14, -3.563) -| (14, -4.063);
  	\draw (14, -2.063) -| (14, -1.563) |- (13.25, -1.563) -| (13.25, -1.813);
  	\draw (14, -4.063) -- (14.75, -4.063);
  	\node[shape=rectangle, minimum width=0.715cm, minimum height=0.465cm](N11) at (14.875, -6.313){} node[anchor=center] at (N11.text){$V_{x}$};
  	\node[shape=rectangle, minimum width=0.965cm, minimum height=0.465cm](N12) at (15.25, -4.063){} node[anchor=center] at (N12.text){$V_{out}$};
  	\node[ground] at (9.75, -8.563){};
  	\node[ground] at (11.25, -8.563){};
  	\node[ground] at (7.5, -1){};
  	\node[ground] at (13.25, -1.813){};
  	\node[shape=rectangle, minimum width=0.965cm, minimum height=0.465cm](N13) at (10, -6.375){} node[anchor=center] at (N13.text){$I_{G_1}$};
  	\node[shape=rectangle, minimum width=4.215cm, minimum height=1.09cm] at (10.25, -0.937){} node[anchor=north, align=center, text width=3.827cm, inner sep=6pt] at (10.25, -0.375){Small Signal Schematic\\(Sink)};
  	\draw (6, 3.52) -| (6, 3.71);
  	\node[nmos, arrowmos](N14) at (6, 4.48){} node[anchor=west] at (N14.text){$M_2$};
  	\draw (15, 6.398) |- (14.77, 6.25);
  	\node[shape=rectangle, minimum width=0.965cm, minimum height=0.465cm](N15) at (2.5, 3.71){} node[anchor=center] at (N15.text){$V_{y}$};
  	\draw (3, 3.71) |- (3.5, 3.71);
  	\node[pmos, arrowmos](N16) at (15, 5.48){} node[anchor=west] at (N16.text){$M_2$};
  	\node[shape=rectangle, minimum width=0.965cm, minimum height=0.465cm](N17) at (11.5, 6.25){} node[anchor=center] at (N17.text){$V_{y}$};
  	\draw (12, 6.25) |- (12.5, 6.25);
  	\draw (12.5, -4.063) to[american current source, l_={$g_{m_2}v_{gs}$}] (12.5, -5.563);
  	\draw (14, -4.063) to[european resistor, l_={$g_{o_2}$}] (14, -5.563);
  	\draw (12.5, -5.563) -- (14, -5.563);
  	\draw (14, -4.063) -- (12.5, -4.063);
  	\draw (12.5, -5.563) -- (11.25, -5.563);
  	\draw (11.25, -4.063) -- (10.5, -4.063);
  	\node[shape=rectangle, minimum width=0.465cm, minimum height=0.465cm] at (11.25, -4.313){} node[anchor=center, align=center, text width=0.077cm, inner sep=6pt] at (11.25, -4.313){\footnotesize G};
  	\node[shape=rectangle, minimum width=0.465cm, minimum height=0.465cm] at (11.25, -5.313){} node[anchor=center, align=center, text width=0.077cm, inner sep=6pt] at (11.25, -5.313){\footnotesize S};
  	\draw (14, -5.563) -- (14, -7.063);
  	\draw (10.5, -7.063) -| (10.5, -6.313) -- (14, -6.313);
  	\draw (14, -6.313) -- (14.5, -6.313);
  	\node[shape=rectangle, minimum width=0.715cm, minimum height=0.465cm](N18) at (5.875, -6.25){} node[anchor=center] at (N18.text){$V_{y}$};
  	\draw (6.25, -6.25) -- (6.75, -6.25);
  	\node[nmos, arrowmos, xscale=-1](N19) at (3.5, 4.48){} node[anchor=east] at (N19.text){$M_3$};
  	\draw (3.5, 3.71) -| (3.5, 3.52);
  	\draw (4.48, 4.48) -- (5.02, 4.48);
  	\draw (4.75, 4.5) |- (3.5, 5.5);
  	\draw (3.5, 5.25) -| (3.5, 5.5);
  	\draw (3.5, 6.04) -| (3.5, 5.5);
  	\node[shape=rectangle, minimum width=0.965cm, minimum height=0.465cm](N20) at (2.5, 5.5){} node[anchor=center] at (N20.text){$V_{z}$};
  	\draw (3, 5.5) |- (3.5, 5.5);
  	\node[shape=rectangle, minimum width=0.965cm, minimum height=0.465cm](N21) at (4.75, 5.75){} node[anchor=center] at (N21.text){$I_{G_2}$};
  	\node[pmos, arrowmos, xscale=-1](N22) at (12.5, 5.48){} node[anchor=east] at (N22.text){$M_3$};
  	\draw (13.48, 5.48) -- (14.02, 5.48);
  	\draw (12.5, 4.562) -- (12.5, 4.71);
  	\draw (12.5, 6.25) |- (12.5, 6.398);
  	\draw (13.75, 5.5) |- (12.5, 4.562);
  	\node[shape=rectangle, minimum width=0.965cm, minimum height=0.465cm](N23) at (13.75, 4.312){} node[anchor=center] at (N23.text){$I_{G_2}$};
  	\draw (10.5, -4.063) -- (9.75, -4.063);
  	\draw (8.25, -4.063) to[american current source, l={$g_{m_3}v_{gs}$}] (8.25, -5.563);
  	\draw (6.75, -4.063) to[european resistor, l={$g_{o_3}$}] (6.75, -5.563);
  	\draw (8.25, -4.063) -- (6.75, -4.063);
  	\draw (8.25, -5.563) -- (6.75, -5.563);
  	\draw (9.75, -5.563) -- (8.25, -5.563);
  	\node[shape=rectangle, minimum width=0.465cm, minimum height=0.465cm] at (9.75, -4.313){} node[anchor=center, align=center, text width=0.077cm, inner sep=6pt] at (9.75, -4.313){\footnotesize G};
  	\node[shape=rectangle, minimum width=0.465cm, minimum height=0.465cm] at (9.75, -5.313){} node[anchor=center, align=center, text width=0.077cm, inner sep=6pt] at (9.75, -5.313){\footnotesize S};
  	\node[shape=rectangle, minimum width=0.465cm, minimum height=0.465cm] at (6.25, -4.313){} node[anchor=center, align=center, text width=0.077cm, inner sep=6pt] at (6.25, -4.313){\footnotesize D};
  	\draw (10.5, -4.063) |- (6.75, -3.25) |- (6.75, -4.063);
  	\node[shape=rectangle, minimum width=0.965cm, minimum height=0.465cm](N24) at (11.5, 4.562){} node[anchor=center] at (N24.text){$V_{z}$};
  	\draw (12, 4.562) |- (12.5, 4.562);
  	\node[shape=rectangle, minimum width=0.715cm, minimum height=0.465cm](N25) at (5.875, -3.25){} node[anchor=center] at (N25.text){$V_{z}$};
  	\draw (6.25, -3.25) -- (6.75, -3.25);
  	\node[shape=rectangle, minimum width=0.965cm, minimum height=0.465cm](N26) at (10.5, -3){} node[anchor=center] at (N26.text){$I_{G_2}$};
  \end{tikzpicture}
\end{center}

\subsubsection{Transfer Function}
Applying KCL at $V_{out}$ on Full Schematic (for Sink):
\begin{align*}
  (V_{out} - V_{cc})G_5 + I_{M_2} &= 0 \\
  V_{out}G_5 - V_{cc}G_5 + I_{M_2} &= 0 \\
  V_{out} &= \frac{V_{cc}G_5 - I_{M_2}}{G_5} \tag{i}
\end{align*}
Applying KCL at $V_{x}$ on Full Schematic (for Sink):
\begin{align*}
  I_{M_1} + I_{G_1} - I_{M_2} &= 0 \\
  I_{M_1} + I_{G_1} &= I_{M_2} \\
  I_{M_1} &\approx I_{M_2} \quad \text{[assuming } I_{G_1} \approx 0 \text{ by MOS physics]} \tag{ii}
\end{align*}
Applying KCL at $V_{y}$ on Full Schematic (for Sink):
\begin{align*}
  I_{M_0} - I_{M_3} &= 0 \\
  I_{M_0} &= I_{M_3} \tag{iii}
\end{align*}
Applying KCL at $V_{z}$ on Full Schematic (for Sink):
\begin{align*}
  (V_{z} - V_{cc})G_4 + I_{M_3} + I_{G_2} &= 0 \\
  V_{z}G_4 - V_{cc}G_4 + (I_{M_3} + I_{G_2}) &= 0 \\
  V_{z}G_4 &= V_{cc}G_4 - (I_{M_3} + I_{G_2}) \\
  V_{z} &= V_{cc} - \frac{(I_{M_3} + I_{G_2})}{G_4} \tag{iv}
\end{align*}
Here's a table of the voltages at each terminal of the
transistors:\\
\begin{tabularx}{\linewidth}{|X|X|X|X|} \hline
  Transistor & Gate (G) & Source (S) & Drain (D) \\ \hline
  $M_0$ & $V_{x}$ & $V_{ee}$ & $V_{y}$ \\ \hline
  $M_1$ & $V_{x}$ & $V_{ee}$ & $V_{x}$ \\ \hline
  $M_2$ & $V_{z}$ & $V_{x}$ & $V_{out}$ \\ \hline
  $M_3$ & $V_{z}$ & $V_{y}$ & $V_{z}$ \\ \hline
\end{tabularx}\\
Here's a table of $V_{gs}$ and $V_{ds}$ for each transistor:\\
\begin{tabularx}{\linewidth}{|X|X|X|} \hline
  Transistor & $V_{gs}$ & $V_{ds}$ \\ \hline
  $M_0$ & $V_{x} - V_{ee}$ & $V_{y} - V_{ee}$ \\ \hline
  $M_1$ & $V_{x} - V_{ee}$ & $V_{x} - V_{ee}$ \\ \hline
  $M_2$ & $V_{z} - V_{x}$ & $V_{out} - V_{x}$ \\ \hline
  $M_3$ & $V_{z} - V_{y}$ & $V_{z} - V_{y}$ \\ \hline
\end{tabularx}\\
Using the above tables, we can express the currents through
each transistor as follows:
\begin{align*}
  I_{M_0} &= k_{n_0}(V_{gs_0} - V_{th})^2 \cdot (1 + \lambda V_{ds_0}) \\
          &= k_{n_0}((V_{x} - V_{ee}) - V_{th})^2 \cdot (1 + \lambda (V_{y} - V_{ee})) \\
          &= k_{n_0}(V_{x} - V_{ee} - V_{th})^2 \quad \text{[assuming } \lambda = 0 \text{]} \tag{v} \\ \\
  I_{M_1} &= k_{n_1}(V_{gs_1} - V_{th})^2 \cdot (1 + \lambda V_{ds_1}) \\
          &= k_{n_1}((V_{x} - V_{ee}) - V_{th})^2 \cdot (1 + \lambda (V_{x} - V_{ee})) \\
          &= k_{n_1}(V_{x} - V_{ee} - V_{th})^2 \quad \text{[assuming } \lambda = 0 \text{]} \tag{vi} \\ \\
  I_{M_2} &= k_{n_2}(V_{gs_2} - V_{th})^2 \cdot (1 + \lambda V_{ds_2}) \\
          &= k_{n_2}((V_{z} - V_{x}) - V_{th})^2 \cdot (1 + \lambda (V_{out} - V_{x})) \\
          &= k_{n_2}\left(\left(V_{cc} - \frac{(I_{M_3} + I_{G_2})}{G_4}\right) - V_{x} - V_{th}\right)^2 \cdot \left(1 + \lambda\left(\frac{V_{cc}G_5 - I_{M_2}}{G_5}\right) - V_{x}\right) \\
          &= k_{n_2}\left(\left(V_{cc} - \frac{(I_{M_3} + I_{G_2})}{G_4}\right) - V_{x} - V_{th}\right)^2 \quad \text{[assuming } \lambda = 0 \text{]} \tag{vii} \\ \\
  I_{M_3} &= k_{n_3}(V_{gs_3} - V_{th})^2 \cdot (1 + \lambda V_{ds_3}) \\
          &= k_{n_3}((V_{z} - V_{y}) - V_{th})^2 \cdot (1 + \lambda (V_{z} - V_{y})) \\
          &= k_{n_3}\left(\left(V_{cc} - \frac{(I_{M_3} + I_{G_2})}{G_4}\right) - V_{y} - V_{th}\right)^2 \codt \left(1 + \lambda\left(\left(V_{cc} - \frac{(I_{M_3} + I_{G_2})}{G_4}\right) - V_{y}\right)\right) \\
          &= k_{n_3}\left(\left(V_{cc} - \frac{(I_{M_3} + I_{G_2})}{G_4}\right) - V_{y} - V_{th}\right)^2 \quad \text{[assuming } \lambda = 0 \text{]} \tag{viii}
\end{align*}
It can be seen that,
\begin{align*}
  I_{in} &= I_{M_3} = I_{M_0} \quad \text{[from (iii)]} \\
  I_{out} &= I_{M_2} = I_{M_1} \quad \text{[from (ii)]} \\
  I_{M_0} &= I_{M_1} \quad \text{[from (v) and (vi) and assuming } k_{n_0} = k_{n_1} \text{]} \\
  \therefore  I_{out} &= I_{in}
\end{align*}
Thus,
\begin{align*}
  \boxed{\frac{I_{out}}{I_{in}} \approx \frac{I_{M_2}}{I_{M_3}} \approx \frac{I_{M_1}}{I_{M_0}} \approx \frac{k_{n_1}}{k_{n_0}} = \frac{\frac{\mu_n C_{ox} W_1}{L_1}}{\frac{\mu_n C_{ox} W_0}{L_0}} = \frac{\frac{W_1}{L_1}}{\frac{W_0}{L_0}} = \frac{W_1 L_0}{L_1 W_0}}
\end{align*}
The analysis is exactly the same for the Source
configuration, with the only difference being the
use of PMOS transistors instead of NMOS transistors.

\subsubsection{Analysis}
\begin{enumerate}
  \item \textbf{Setting Current to be Mirrored}:
    \begin{enumerate}
      \item Choose appropriate transistor dimensions
        ($W$ and $L$) for transistors $M_0$ and $M_1$
        to achieve the desired current gain.
        \begin{align*}
          \frac{I_{out}}{I_{in}} &= \frac{W_1 L_0}{L_1 W_0}
        \end{align*}
      \item Choose a suitable overdrive voltage value
        ($V_{ov}$) for computing the bias voltages at
        nodes $V_{x}$, $V_{y}$, and $V_{z}$.
        \begin{align*}
          V_{ov_0} &= V_{gs_0} - V_{th} \\
          V_{ov_0} &= V_{x} - V_{ee} - V_{th} \\
          V_{x} &= V_{ov_0} + V_{ee} + V_{th} \\
          \\
          V_{ov_2} &= V_{gs_2} - V_{th} \\
          V_{ov_2} &= V_{z} - V_{x} - V_{th} \\
          V_{z} &= V_{ov_2} + V_{x} + V_{th} \\
          \\
          V_{ov_3} &= V_{gs_3} - V_{th} \\
          V_{ov_3} &= V_{z} - V_{y} - V_{th} \\
          V_{y} &= V_{z} - V_{ov_3} - V_{th}
        \end{align*}
      \item Using the chosen $V_{ov}$, compute $W_0$
          to set the required bias current $I_{in}$.
          \begin{align*}
            I_{in} = I_{M_0} &= \frac{1}{2} \mu_n C_{ox} \frac{W_0}{L_0} V_{ov_0}^2 \\
            \therefore W_0 &= \frac{2 I_{in} L_0}{\mu_n C_{ox} V_{ov_0}^2}
          \end{align*}
          $W_1$ can then be computed using the current
          gain relation from step (a).
      \item Compute the required value of $G_4$ (or
          $R_4$) to ensure that the voltage at node
          $V_{z}$ is as calculated in step (b).
          \begin{align*}
            V_{z} &= V_{cc} - \frac{(I_{M_3} + I_{G_2})}{G_4} \\
            V_{z} &\approx V_{cc} - \frac{I_{M_3}}{G_4} \quad \text{[assuming } I_{G_2} \approx 0 \text{ by MOS physics]} \\
            V_{z} - V_{cc} &\approx - \frac{I_{M_3}}{G_4} \\
            \frac{(I_{M_3})}{(V_{cc} - V_{z})} &\approx G_4 \\
            \frac{(I_{in})}{(V_{cc} - V_{z})} &\approx G_4 \quad \text{[} I_{in} = I_{M_3} = I_{M_0} \text{]} \\
            \therefore R_4 &\approx \frac{V_{cc} - V_{z}}{I_{in}}
          \end{align*}
    \end{enumerate}
  \item \textbf{Output Load Limit}: Maximum possible
    load that can be connected at the output before
    the transistor $M_2$ goes out of saturation can
    be calculated as follows:
    \begin{align*}
      R_{out, max} &= \frac{V_{L, max}}{I_{M_2}} \quad \text{[where, } V_{ee} \leq V_{L, max} \leq V_{cc} \text{]}
    \end{align*}
  \item \textbf{Impedances}:
    \paragraph{At Input:} \textbf{Ignoring $G_4$
      which is used to set the current,} the input
      impedance is simply the output resistance of
      transistor $M_0$ scaled by the gain of
      transistor $M_3$.
      \begin{align*}
        Z_{in} &\approx \frac{1}{g_{o_0}\frac{g_{o_3}}{g_{m_3}}} \\
        &\approx \frac{1}{(\lambda_{M_0} I_{M_0}) \cdot \frac{(\lambda_{M_3} I_{M_3})}{\left(\frac{2 I_{M_3}}{V_{ov_3}}\right)}} \\
        &\approx \frac{1}{\lambda_{M_0} I_{M_0}} \cdot \frac{1}{\frac{\lambda_{M_3} V_{ov_3}}{2}} \\
        &\approx \frac{1}{\lambda_{M_0} I_{M_0}} \cdot \frac{2}{\lambda_{M_3} V_{ov_3}}\\
        &\approx \frac{2}{\lambda_{M_0} \lambda_{M_3} I_{M_0} (V_{gs_3} - V_{th})} \\
        &\approx \frac{2}{\lambda_{M_0} \lambda_{M_3} I_{M_0} ((V_{z} - V_{y}) - V_{th})}
      \end{align*}
    \paragraph{At Output:} \textbf{Ignoring $G_5$
      which is used as a load,} the output impedance
      is the output resistance of transistor $M_1$
      scaled by the gain of transistor $M_2$.
      \begin{align*}
        Z_{out} &\approx \frac{1}{g_{o_1}\frac{g_{o_2}}{g_{m_2}}} \\
        &\approx \frac{1}{(\lambda_{M_1} I_{M_1}) \cdot \frac{(\lambda_{M_2} I_{M_2})}{\left(\frac{2 I_{M_2}}{V_{ov_2}}\right)}} \\
        &\approx \frac{1}{\lambda_{M_1} I_{M_1}} \cdot \frac{1}{\frac{\lambda_{M_2} V_{ov_2}}{2}} \\
        &\approx \frac{1}{\lambda_{M_1} I_{M_1}} \cdot \frac{2}{\lambda_{M_2} V_{ov_2}}\\
        &\approx \frac{2}{\lambda_{M_1} \lambda_{M_2} I_{M_1} (V_{gs_2} - V_{th})} \\
        &\approx \frac{2}{\lambda_{M_1} \lambda_{M_2} I_{M_1} ((V_{z} - V_{x}) - V_{th})}
      \end{align*}
\end{enumerate}

\subsubsection{Example Design}

\paragraph{Question:}
Design a Modified Wilson Current Mirror in
Sink configuration to achieve a current gain of
2, with an input bias current of $I_{in} =
10\mu A$. Assume the following parameters for
the NMOS transistors used:
\begin{align*}
  \mu_n C_{ox} = 100 \mu A/V^2 \spc V_{th} = 0.7 V \spc \lambda = 0.02 V^{-1} \\
  V_{cc} = 5 V \scp V_{ee} = 0 V \spc L_0 = L_1 = L_2 = L_3 = 1 \mu m
\end{align*}
Find the output impedance and the maximum
possible load that can be connected at the output
assuming $V_{L,max} = V_{cc}$
\paragraph{Solution:}
\begin{enumerate}
  \item \textbf{Choose appropriate transistor dimensions:}
    \begin{align*}
      \frac{I_{out}}{I_{in}} &= \frac{W_1 L_0}{L_1 W_0} \\
      2 &= \frac{W_1 \cdot 1 \mu m}{1 \mu m \cdot W_0} \\
      W_1 &= 2 \cdot W_0
    \end{align*}
  \item \textbf{Choose a suitable overdrive voltage and determine
    the resulting bias voltages:}
    \begin{align*}
      V_{ov_0} &= 0.2 V \\
      V_{x} &= V_{ov_0} + V_{ee} + V_{th} = 0.2 + 0 + 0.7 = 0.9 V \\
      \\
      V_{ov_2} &= 0.2 V \\
      V_{z} &= V_{ov_2} + V_{x} + V_{th} = 0.2 + 0.9 + 0.7 = 1.8 V \\
      \\
      V_{ov_3} &= 0.2 V \\
      V_{y} &= V_{z} - V_{ov_3} - V_{th} = 1.8 - 0.2 - 0.7 = 0.9 V
    \end{align*}
  \item \textbf{Compute $W_0$ to set the required bias current:} For $I_{in} = 10 \mu A$
    \begin{align*}
      W_0 &= \frac{2 I_{in} L_0}{\mu_n C_{ox} V_{ov_0}^2} \\
          &= \frac{2 \cdot 10 \mu A \cdot 1 \mu m}{100 \mu A \cdot (0.2)^2} \\
          &= 5 \mu m
    \end{align*}
    Using (i) and (ii):
    \begin{align*}
      W_1 &= 2 \cdot W_0 = 2 \cdot 5 \mu m = 10 \mu m
    \end{align*}
  \item \textbf{Compute the required value of $G_4$ (or $R_4$):}
    \begin{align*}
      R_4 &\approx \frac{V_{cc} - V_{z}}{I_{in}} \\
          &= \frac{5 V - 1.8 V}{10 \mu A} \\
          &= 320 k\Omega
    \end{align*}
  \item \textbf{Output Load Limit:}
    \begin{align*}
      R_{out, max} &= \frac{V_{L, max}}{I_{M_2}} = \frac{5 V}{10 \mu A} = 500 k\Omega
    \end{align*}
  \item \textbf{Impedances:}
    \paragraph{At Input:}
    \begin{align*}
      Z_{in} &\approx \frac{2}{\lambda_{M_0} \lambda_{M_3} I_{M_0} (V_{gs_3} - V_{th})} \\
             &= \frac{2}{0.02 \imes 0.02 \times 10 \mu A \times (1.8 - 0.9)} \\
             &= 5.56 G\Omega
    \end{align*}
    \paragraph{At Output:}
    \begin{align*}
      Z_{out} &\approx \frac{2}{\lambda_{M_1} \lambda_{M_2} I_{M_1} (V_{gs_2} - V_{th})} \\
              &= \frac{2}{0.02 \imes 0.02 \times 20 \mu A \times (1.8 - 0.9)} \\
              &= 2.78 G\Omega
    \end{align*}
\end{enumerate}
